% Generated draft: requires mathematical and source review.
\documentclass[11pt]{article}
\usepackage[margin=25mm]{geometry}
\usepackage{fontspec}
\setmainfont{Latin Modern Roman}
\usepackage{amsmath,amssymb,mathtools,unicode-math}
\setmathfont{Latin Modern Math}
\usepackage{enumitem,needspace,xurl,hyperref}
\setlength{\parindent}{0pt}
\setlength{\parskip}{5pt}
\setlength{\emergencystretch}{3em}
\newcommand{\N}{\mathbb{N}}
\newcommand{\Z}{\mathbb{Z}}
\newcommand{\Q}{\mathbb{Q}}
\newcommand{\R}{\mathbb{R}}
\newcommand{\C}{\mathbb{C}}
\newcommand{\F}{\mathbb{F}}
\newcommand{\A}{\mathbb{A}}
\newcommand{\PP}{\mathbb{P}}
\newcommand{\E}{\mathbb{E}}
\newcommand{\eps}{\varepsilon}
\newcommand{\dd}{\,\mathrm{d}}
\begin{document}
\section*{UnsolvedMath ID: 30004834}
\textbf{Source problem code:} catalogue entry
\subsection{Definitions and mathematical statement}
The supplied record does not determine a unique mathematical problem. Its clipped statement contains only ``question: [1, Question 6.8]'' and a reference to page 2425. The supplied literature assessment interprets this as a question about uniqueness of a \(C^*\)-norm on a group algebra, whereas the supplied statement-verification note says that the cited construction settles Question 6.8. The title instead refers to a resolved Witt-group question about anisotropic modular categories. These descriptions cannot be identified with one another from the supplied evidence.

To make the group-algebra interpretation explicit, let \(G\) be an abstract group and interpret its group algebra as the complex algebra
\[
\C[G]=\left\{\sum_{g\in G}a_g g : a_g\in\C,\ a_g=0\text{ for all but finitely many }g\right\},
\]
with multiplication induced by the group law and involution
\[
\left(\sum_{g\in G}a_g g\right)^*=\sum_{g\in G}\overline{a_g}\,g^{-1}.
\]
A \(C^*\)-norm on this algebra means a norm, rather than merely a seminorm, satisfying \(\|ab\|\leq\|a\|\|b\|\), \(\|a^*\|=\|a\|\), and \(\|a^*a\|=\|a\|^2\). Uniqueness means that any two such norms agree on every element of \(\C[G]\). The group \(G\) is locally finite if every subgroup generated by finitely many elements is finite.

For a specified class \(\mathcal G\) of groups, the question suggested by the supplied literature assessment would be whether
\[
\forall G\in\mathcal G,\qquad
\bigl(\C[G]\text{ admits exactly one }C^*\text{-norm}\bigr)
\Longrightarrow
\bigl(\forall F\subseteq G\text{ finite},\ |\langle F\rangle|<\infty\bigr).
\]
This is only a conditional reconstruction. The class \(\mathcal G\), any additional hypotheses, and the identification of this implication with the intended source problem remain unresolved. The supplied record does not provide enough information to formulate the alternative Witt-group target.
\subsection{Short English statement}
The supplied assessment suggests asking whether a unique \(C^*\)-norm on a complex group algebra forces every finitely generated subgroup to be finite. However, the fragment, title, and supplied verification note conflict, and the intended class of groups is unspecified. The displayed question is therefore a provisional interpretation requiring mathematical and source review.
\subsection{Sources}
\begin{itemize}
\item Supplied local registry record, permanent UnsolvedMath ID \(30004834\), exact source problem code \texttt{catalogue entry}. The record labels the problem open but marks its statement \texttt{unrecoverable}; no proposer is supplied.
\item Supplied report citation: \emph{Quantum Groups -- Algebra, Analysis and Category Theory} (2022), Oberwolfach Reports, DOI \texttt{10.4171/owr/2021/44}, pp.~2397--2458. The clipped statement points to p.~2425 and ``[1, Question 6.8]''. Supplied literature-source label: ``Oberwolfach Report 44''. \url{https://doi.org/10.4171/owr/2021/44}.
\item Supplied literature source: \emph{Just-infinite C-star algebras}. \url{https://ems.press/content/serial-article-files/43460}. The supplied assessment associates Question 6.8 with the group-algebra implication above; the paper's text and bibliographic authorship are not included in the record.
\item Literature assessment dated exactly \texttt{2026-08-22}: the supplied curation labels the group-algebra converse open and says the published source records it as unresolved. This is supplied dated curation, not independently verified. The background expressly requests independent verification of time-sensitive status and cited claims.
\item Conflicting supplied provenance: the statement-verification note reports that the linked report was inspected and that its construction settles the cited Question 6.8. The title is ``Resolved Witt-Group Question on Anisotropic Modular Categories''. These are supplied descriptions, not independent verification performed for this draft. The record gives proposed year 2021, while its report citation gives 2022; both are retained without reconciliation.
\end{itemize}
\end{document}
